\section{A positive gamma action and its integer refinement}\label{app:gamma}
This appendix records the common positive bulk from which a different
joint source could be built. For $\alpha>0$, let
\begin{equation}\label{eq:alpha-tower}
 a_n(\alpha)=2(n+\alpha),\qquad
 b_\alpha(s)=\sum_{n\geq0}\frac2{a_n(\alpha)}
                   \frac{s}{a_n(\alpha)^2+s}.
\end{equation}
Thus $b=b_{1/4}$. For line fields $u_n$ with weak derivatives, define
the nonnegative joint energy
\begin{equation}\label{eq:alpha-action}
 \mathcal E_\alpha(F,\{u_n\})
 =\sum_{n\geq0}\frac2{a_n(\alpha)}
 \left(\norm{F-u_n}^2+a_n(\alpha)^{-2}\norm{u_n'}^2\right).
\end{equation}
The finite-energy domain is specified by the sum as a whole.
In particular, the divergent expression
$\sum_n2\norm F^2/a_n(\alpha)$ must not be split off as a separate
energy from compensating fields.

\begin{proposition}[Elimination and positive refinement]\label{prop:refinement}
For smooth compactly supported $F$,
\begin{equation}\label{eq:alpha-inf}
 \inf_{\{u_n\}}\mathcal E_\alpha(F,\{u_n\})
 =\frac1{2\pi}\int_\R b_\alpha(\tau^2)|\widehat F(\tau)|^2\dd\tau.
\end{equation}
For each integer $p\geq2$, put $\alpha_j=(\alpha+j)/p$,
$0\leq j<p$, and $D_pF(x)=p^{-1/2}F(x/p)$. Then
\begin{align}
 b_\alpha(s)&=\frac1p\sum_{j=0}^{p-1}b_{\alpha_j}(s/p^2),
 \label{eq:refine-multiplier}\\
 \mathcal E_\alpha(F,\{u_n\})
 &=\frac1p\sum_{j=0}^{p-1}
   \mathcal E_{\alpha_j}
   (D_pF,\{D_pu_{pm+j}\}_{m\geq0}).
 \label{eq:refine-action}
\end{align}
The latter identity is valid as an equality of nonnegative extended
energies, before eliminating the fields.
\end{proposition}
\begin{proof}
For each mode, minimizing the Fourier-space quadratic function gives
$\widehat u_n=a_n^2(a_n^2+\tau^2)^{-1}\widehat F$ and the summand
in \eqref{eq:alpha-tower}. Substitution of all minimizers, or finite
towers followed by monotone convergence, proves \eqref{eq:alpha-inf}.
For refinement, write $n=pm+j$. Then
$a_{pm+j}(\alpha)=p\,a_m(\alpha_j)$,
$\norm{D_pu}=\norm u$, and
$\norm{(D_pu)'}=p^{-1}\norm{u'}$. Substitution proves
\eqref{eq:refine-action}; reindexing the positive multiplier sum gives
\eqref{eq:refine-multiplier}.
\end{proof}

The first examples are
\begin{align}
 b_{1/4}(s)&=\tfrac12b_{1/8}(s/4)+\tfrac12b_{5/8}(s/4),
 \notag\\
 b_{1/4}(s)&=\tfrac13\left[b_{1/12}(s/9)+b_{5/12}(s/9)
                           +b_{3/4}(s/9)\right].
 \label{eq:refine-examples}
\end{align}
Refining by $p$ and then by another integer $\ell$ gives residue
labels $j+pk$, the dilation $D_{p\ell}$ and weights $1/(p\ell)$.
Thus refinement by $2$ then $3$, or $3$ then $2$, gives the same six
branches $\{1,5,9,13,17,21\}/24$ up to permutation.

There is also a contact identity. Set
$m_\alpha(\tau)=\re\psi(\alpha+i\tau/2)-\log\pi$. Gauss
multiplication \cite{DLMF} gives
\begin{equation}\label{eq:refine-contact}
 m_\alpha(\tau)=\log p+
       \frac1p\sum_{j=0}^{p-1}m_{\alpha_j}(\tau/p).
\end{equation}
Subtracting the value at zero recovers \eqref{eq:refine-multiplier}.
Only the kinetic actions in \eqref{eq:alpha-action} have the asserted
positivity; \eqref{eq:refine-contact} is not a positive decomposition
of the complete archimedean form.

These formulas exist for every integer, and so do not select primes.
Coordinate dilation by $p$ is also different from translation by
$\log p$. A conservative junction acting on the common refined fields
could be tested, but its injection, return and readout maps would have
to derive these relations. The normalization $p^{-1/2}$ in the
residual-output inclusion is fixed by \eqref{eq:refine-action}; it is
not yet a physical arithmetic return amplitude. Refinement or a unitary
rotation of all outputs preserves their norm and cannot itself introduce
prime terms. No novelty claim is made for gamma multiplication or its
semilocal use \cite{QuasiInner}.
